% from ega4-2-fr.tex
\oldpage[IV-2]{176}
V) \emph{Fin de la démonstration.} --- D'après II), on peut supposer que
$B$ est une $A$-algèbre intègre de type fini contenant $A$. Soient $K$ le
corps des fractions de $A$, $L$ celui de $B$, qui est une extension de type
fini de $K$. Soit $M$ la fermeture algébrique de $K$ dans $L$, qui est une
extension algébrique finie de $K$; soit $C_0$ la fermeture intégrale de $A$
dans $M$, qui est aussi la fermeture intégrale de $A$ dans $L$; on a donc
$C=B\cap C_0$; si l'on pose $C'_0=C_0\otimes_A A'$, on a par suite
$C'=B'\cap C'_0$ par platitude (0$_{\mathrm I}$, 6.1.3). Or, il résulte de
IV) que $C'_0$ est la fermeture intégrale de $A'$ dans
$M'=M\otimes_A A'$; en outre, $M'$ est un anneau noethérien et le morphisme
$\operatorname{Spec}(M')\to\operatorname{Spec}(M)$ est normal (comme on
l'a vu dans IV)); appliquant III) à $M$ et $M'$ au lieu de $A$ et $A'$ et
à $L$ au lieu de $B$, on en déduit que $M'$ est intégralement fermé dans
$L'=L\otimes_A A'$; donc $C'_0$ est la fermeture intégrale de $A'$ dans
$L'$, et $C'=C'_0\cap B'$ est la fermeture intégrale de $A'$ dans $B'$.
C.Q.F.D.
\end{proof}

\begin{corollary}[6.14.5]
\label{IV.2.s6.14.5-fr}
Soient $A$ un anneau noethérien, $A'$ une $A$-algèbre noethérienne telle
que le morphisme $\operatorname{Spec}(A')\to\operatorname{Spec}(A)$ soit
normal. Soient $B$, $C$ deux $A$-algèbres,
$\varphi:B\to C$ un $A$-homomorphisme faisant de $C$ une $B$-algèbre, $D$
la fermeture intégrale de $B$ dans $C$. Si l'on pose
$B'=B\otimes_A A'$, $C'=C\otimes_A A'$, $D'=D\otimes_A A'$, $D'$ est la
fermeture intégrale de $B'$ dans $C'$.
\end{corollary}

\begin{proof}
Soit $(B_\lambda)$ la famille filtrante croissante des sous-$A$-algèbres
de type fini de $B$, et posons $B'_\lambda=B_\lambda\otimes_A A'$, de
sorte que $B'$ est réunion de la famille filtrante croissante des
sous-$A'$-algèbres de type fini $B'_\lambda$. Soit $D_\lambda$ la
fermeture intégrale de $B_\lambda$ dans $C$, et posons
$D'_\lambda=D_\lambda\otimes_A A'$, de sorte que $D$ est réunion des
$D_\lambda$ et $D'$ réunion des $D'_\lambda$; si nous prouvons que
$D'_\lambda$ est la fermeture intégrale de $B'_\lambda$ dans $C'$, il en
résultera que $D'$ est la fermeture intégrale de $B'$ dans $C'$. Or
$B_\lambda$ et $B'_\lambda$ sont noethériens et le morphisme
$\operatorname{Spec}(B'_\lambda)\to\operatorname{Spec}(B_\lambda)$ est
normal (6.8.2); on peut donc appliquer (6.14.4) en remplaçant $A$, $A'$ et
$B$ par $B_\lambda$, $B'_\lambda$ et $C$ respectivement, d'où le
corollaire.
\end{proof}

On peut par exemple appliquer (6.14.5) lorsque $A$ est un anneau local
excellent et $A'$ son complété $\widehat A$, puisque dans ce cas
$\operatorname{Spec}(A')\to\operatorname{Spec}(A)$ est un morphisme
régulier (7.8.2).

\subsection{Préschémas géométriquement unibranches}
\label{IV.2.s6.15-fr}

\begin{env}[6.15.1]
\label{IV.2.s6.15.1-fr}
Nous dirons qu'un préschéma $X$ est \emph{unibranche} (resp.
\emph{géométriquement unibranche}) en un point $x$, ou que le point $x$
est unibranche (resp. géométriquement unibranche) dans $X$ si l'anneau
local $\mathcal O_{X,x}$ est unibranche (resp. géométriquement
unibranche) (0, 23.2.1); nous dirons que $X$ est unibranche (resp.
géométriquement unibranche) s'il l'est en tout point. Il suit de cette
définition que, pour que $X$ soit unibranche (resp. géométriquement
unibranche) en un point, il faut et il suffit que $X_{\mathrm{red}}$ le soit
en ce point.
\end{env}

\begin{env}[6.15.2]
\EGAReaderTarget{ega4.c04.statement.6-15-2}{IV.2.s6.15.2-fr}
Il faut prendre garde qu'avec les définitions de (6.15.1), un anneau local
$A$ peut être unibranche (resp. géométriquement unibranche) sans que
$\operatorname{Spec}(A)$ le soit; en d'autres termes, il peut se faire qu'il
y ait des idéaux premiers $\mathfrak p$ de $A$ tels que $A_\mathfrak p$ ne
soit pas unibranche; il revient aussi au même de dire que sur un préschéma
la notion de point géométriquement unibranche n'est pas stable par
générisation. Nous allons le voir sur l'exemple suivant.
\end{env}

